\section{研究の目的}

本研究の目的は，軟弱地面において床沈み込みが発生する環境下でも安定した二足歩行を実現するための，重心高さ補償型歩行生成アルゴリズムを提案することである．

具体的には，床沈み込みにより生じる支持脚高さの変動を推定し，その変動成分を重心高さの参照値に反映させることで，動的に変化する床高さに追従可能な歩行軌道を生成する．沈み込みが検出された場合には，鉛直方向の速度および加速度項を目標重心高さに加算し，重心の鉛直運動を能動的に制御することで，水平方向の安定性低下を抑制することを目指す．

また，本研究では，LIPMの基本構造を維持した拡張手法を採用することで，計算コストの増大を抑えつつ，既存の歩行制御フレームワークへの統合を容易にすることを重視する．これにより，従来のLIPMベース歩行制御が抱える軟弱地面に対する脆弱性を補完し，実環境適応性の向上を図る．